\chapter{Chords: Notes Played Together}
\label{ch:chords-by-shape}
\chapterline

So far you have played notes one after another. A tune you can hum works that way: one note, then the next. A line of notes heard one after another is called a \textbf{melody}.

Now press three keys at the same moment. The notes blend into one fuller sound. Three or more notes played together make a \textbf{chord}. The sound of chords supporting a melody is called \textbf{harmony}.

Think of a singer in a room. The melody is the voice. The harmony is the lighting: the voice stays the same, but the mood of the room changes. Hear a short melody alone, \listen{E4 D4 C4:2} and then the same melody with chords under it. \listen{C3+E3+G3+E4 G2+B2+D3+D4 C3+E3+G3+C4:2}

In this chapter you will build chords with one hand. You will hear why some chords sound bright and others sound sad. And you will meet the small family of chords that most songs use.

\section{Your first chord: C major}

Put your right thumb on middle C (C4, from Chapter~1). This is where the chord starts, and it gives the chord its name. The note a chord is built on is called its \textbf{root}.

Now count up 4 keys from C. Count every key, black or white, and do not count the key you start on. You pass C\sh, D and D\sh, and land on E. Counting keys like this is counting half steps, as in Chapter~1.

From E, count up 3 more keys: F, F\sh, G. You land on G. Play C, E and G together. This is the \textbf{C major chord}. \listen{C4+E4+G4:2} Its chord symbol is the letter C on its own.

Use fingers 1, 3 and 5: thumb, middle finger and little finger. Fingers 2 and 4 float above the keys they skip. Hear the three notes one at a time, then together. \listen{C4 E4 G4 C4+E4+G4:2}

A chord of three notes built this way is called a \textbf{triad} (``tri'' means three). Each of its notes has a name:
\begin{itemize}
  \item the root is the bottom note, C, which names the chord.
  \item the \textbf{third} is the middle note, E. It is three letters up from the root: C, D, E. (C up to E is the major third from Chapter~1.)
  \item the \textbf{fifth} is the top note, G. It is five letters up: C, D, E, F, G. (C up to G is the fifth from Chapter~1.)
\end{itemize}

\begin{trapbox}
When you count 4 keys, count the black keys too. If you count only white keys, you land on G instead of E. A good check: the major chord on C uses only white keys, and your hand looks like ``play one, skip one, play one, skip one, play one''.
\end{trapbox}

\section{Major and minor: one key makes the difference}

Keep your thumb on C and your little finger on G. Now move only your middle finger one key to the left, from E down to the black key E\fl. Play C, E\fl{} and G together. \listen{C4+Eb4+G4:2}

This is the \textbf{C minor chord}. Its symbol is Cm; the small m means minor. Play C major and then C minor, and listen. \listen{C4+E4+G4:2 C4+Eb4+G4:2} Most people hear major as bright and settled, and minor as darker or sadder.

The root did not change, and neither did the top note. Only the middle note moved, by one key. That gives two short recipes, which work from any root:
\begin{itemize}
  \item Major: from the root, count up 4 keys, then 3 more. We will write this as ``4 then 3''.
  \item Minor: from the root, count up 3 keys, then 4 more: ``3 then 4''.
\end{itemize}
In both, the top note is 7 keys above the root. Figure~\ref{fig:major-vs-minor-triad} shows the two side by side.

\begin{figure}[htbp]
\centering
\begin{tikzpicture}
  \def\wk{0.8}
  \def\hk{2.6}
  \foreach \sy/\ttl/\midx/\aa/\ab in {0/{C major (4 then 3)}/2.5/4/3, -5.0/{C minor (3 then 4)}/2.0/3/4} {
    \begin{scope}[shift={(0,\sy)}]
      \node[font=\small, anchor=south] at ({3.5*\wk},{\hk+0.95}) {\ttl};
      \foreach \k in {0,...,6} {
        \draw[black!60, line width=0.6pt, fill=white] ({\k*\wk},0) rectangle ++(\wk,\hk);
      }
      \draw[black!60, line width=0.6pt, fill=sourceblue!25] (0,0) rectangle ++(\wk,\hk);
      \draw[black!60, line width=0.6pt, fill=sourceblue!25] ({4*\wk},0) rectangle ++(\wk,\hk);
      \foreach \k in {0,1,3,4,5} {
        \fill[black!85] ({(\k+1)*\wk-0.24},{0.38*\hk}) rectangle ++(0.48,{0.62*\hk});
      }
      \node[font=\small] at ({0.5*\wk},0.35) {C};
      \node[font=\small] at ({4.5*\wk},0.35) {G};
      \node[font=\small, text=black!60] at ({0.5*\wk},-0.35) {1};
      \node[font=\small, text=black!60] at ({\midx*\wk},-0.35) {3};
      \node[font=\small, text=black!60] at ({4.5*\wk},-0.35) {5};
      \draw[thick, sourceblue!80!black] ({0.5*\wk},{\hk+0.1}) to[out=60,in=120] node[above, font=\small] {\aa} ({\midx*\wk},{\hk+0.1});
      \draw[thick, sourceblue!80!black] ({\midx*\wk},{\hk+0.1}) to[out=60,in=120] node[above, font=\small] {\ab} ({4.5*\wk},{\hk+0.1});
    \end{scope}
  }
  % middle notes: E (major), E flat (minor)
  \draw[black!60, line width=0.6pt, fill=sourceblue!25] ({2*0.8},0) rectangle ++(0.8,2.6);
  \fill[black!85] ({2*0.8-0.24},{0.38*2.6}) rectangle ++(0.48,{0.62*2.6});
  \fill[black!85] ({3*0.8-0.24},{0.38*2.6}) rectangle ++(0.48,{0.62*2.6});
  \node[font=\small] at ({2.5*0.8},0.35) {E};
  \fill[sourceblue!80!black] ({2*0.8-0.24},{-5.0+0.38*2.6}) rectangle ++(0.48,{0.62*2.6});
  \node[font=\small, text=white] at ({2*0.8},{-5.0+0.38*2.6+0.35}) {E$\flat$};
\end{tikzpicture}
\caption{C major and C minor share C and G. Only the middle note moves, from E down to E\fl. The small numbers under the keys are the fingers to use.}
\label{fig:major-vs-minor-triad}
\end{figure}

Try the recipe on D. Count 4 keys up from D: D\sh, E, F, F\sh. Then 3 more: G, G\sh, A. So D major is D, F\sh{} and A. \listen{D4+F#4+A4:2} Move the middle finger down one key to F, and you have D minor, Dm: D, F and A. \listen{D4+F4+A4:2}

Now A. A major is A, C\sh{} and E, \listen{A3+C#4+E4:2} and A minor, Am, is A, C and E. \listen{A3+C4+E4:2}

Notice that a triad always skips a letter between its notes: D, (E), F\sh, (G), A. That is why we call the middle note of D major F\sh{} and not G\fl, even though it is the same key. As with scales in Chapter~2, the letter name follows the pattern.

\begin{tryitbox}
Build G major and G minor, then E major and E minor, by counting. Play each pair and listen for the change of mood. You should find G, B, D and G, B\fl, D, \listen{G3+B3+D4:2 G3+Bb3+D4:2} then E, G\sh, B and E, G, B. \listen{E4+G#4+B4:2 E4+G4+B4:2}
\end{tryitbox}

\section{A shortcut on the white keys}

Counting works from any key. On the white keys there is a quicker way to find a chord: play a white key, skip a white key, play, skip, play. From C that gives C, E and G.

Try that shape starting on each white key in turn. \listen{C4+E4+G4 D4+F4+A4 E4+G4+B4 F4+A4+C5 G4+B4+D5 A4+C5+E5 B4+D5+F5 C5+E5+G5:2} Some of the chords sound bright and some sound dark:
\begin{itemize}
  \item starting on C, F or G gives a major chord: C (C, E, G), F (F, A, C) and G (G, B, D);
  \item starting on D, E or A gives a minor chord: Dm (D, F, A), Em (E, G, B) and Am (A, C, E);
  \item starting on B gives B, D and F, which is neither. Count it: 3 keys, then 3 more. We meet this chord in the next section.
\end{itemize}

Why does one hand shape give different chords? Because the white keys are not evenly spaced. There is no black key between E and F, or between B and C. So the same ``skip a white key'' gap is sometimes 4 keys and sometimes 3. When in doubt, count.

\section{Two more sounds: sus and diminished}

Play C major, then lift your middle finger off E and put finger 4 on F. You now hold C, F and G. \listen{C4+F4+G4:2} This is \textbf{Csus4} (said ``C sus four''). ``Sus'' is short for suspended: the middle note is hung up one key too high, so the chord is neither major nor minor.

A sus chord sounds open and unfinished. It wants to fall back to the major chord. Hear the pull from Csus4 to C. \listen{C4+F4+G4:2 C4+E4+G4:3} You can also use finger 2 on D instead: C, D and G is \textbf{Csus2}. \listen{C4+D4+G4:2 C4+E4+G4:3}

The chord you found on B, with B, D and F, is a \textbf{diminished} chord. \listen{B3+D4+F4:2} Its recipe is 3 then 3. The top note is only 6 keys above the root instead of 7, and the chord sounds tense. It is written Bdim, or B\textdegree{} with a small circle. Hear it lean into C major. \listen{B3+D4+F4:2 C4+E4+G4:3}

Here are all the recipes so far, with C as the root. You do not need to memorize the table; it is here to look back at.

\begin{center}
\begin{tabular}{llll}
\toprule
Chord & Symbol & Count up & Notes from C \\
\midrule
Major & C & 4 then 3 & C, E, G \\
Minor & Cm & 3 then 4 & C, E\fl, G \\
Diminished & Cdim & 3 then 3 & C, E\fl, G\fl \\
Sus4 & Csus4 & 5 then 2 & C, F, G \\
Sus2 & Csus2 & 2 then 5 & C, D, G \\
\bottomrule
\end{tabular}
\end{center}

\section{Three hand shapes for most major chords}

Counting is reliable, but pianists do not count while they play. They remember how each chord looks and feels under the hand. Nine of the twelve major chords fall into three shapes, shown in Figure~\ref{fig:three-chord-shapes}:
\begin{itemize}
  \item White, white, white: C, F and G major. \listen{C4+E4+G4 F4+A4+C5 G4+B4+D5:2}
  \item White, black, white: D, E and A major. The middle finger sits up on a black key. \listen{D4+F#4+A4 E4+G#4+B4 A3+C#4+E4:2}
  \item Black, white, black: D\fl, E\fl{} and A\fl{} major. Thumb and little finger sit on black keys, with the middle finger down on a white key. \listen{Db4+F4+Ab4 Eb4+G4+Bb4 Ab3+C4+Eb4:2} The notes are D\fl, F, A\fl; E\fl, G, B\fl; and A\fl, C, E\fl.
\end{itemize}

\begin{figure}[htbp]
\centering
\begin{tikzpicture}
  \def\wk{0.58}
  \def\hk{2.3}
  \foreach \sy/\ttl/\sub in {0/{C, F, G}/{white, white, white}, -3.2/{D, E, A}/{white, black, white}, -6.4/{D$\flat$, E$\flat$, A$\flat$}/{black, white, black}} {
    \begin{scope}[shift={(0,\sy)}]
      \node[font=\small, anchor=west, align=left] at ({7*\wk+0.4},{0.5*\hk}) {\ttl\\[2pt]{\color{black!60}\sub}};
      \foreach \k in {0,...,6} {
        \draw[black!60, line width=0.6pt, fill=white] ({\k*\wk},0) rectangle ++(\wk,\hk);
      }
    \end{scope}
  }
  % C E G
  \foreach \k/\n in {0/C,2/E,4/G} {
    \draw[black!60, line width=0.6pt, fill=sourceblue!25] ({\k*\wk},0) rectangle ++(\wk,\hk);
    \node[font=\small] at ({(\k+0.5)*\wk},0.32) {\n};
  }
  % D F# A
  \foreach \k/\n in {1/D,5/A} {
    \draw[black!60, line width=0.6pt, fill=sourceblue!25] ({\k*\wk},-3.2) rectangle ++(\wk,\hk);
    \node[font=\small] at ({(\k+0.5)*\wk},{-3.2+0.32}) {\n};
  }
  % Db F Ab
  \draw[black!60, line width=0.6pt, fill=sourceblue!25] ({3*\wk},-6.4) rectangle ++(\wk,\hk);
  \node[font=\small] at ({3.5*\wk},{-6.4+0.32}) {F};
  \foreach \sy in {0,-3.2,-6.4} {
    \foreach \k in {0,1,3,4,5} {
      \fill[black!85] ({(\k+1)*\wk-0.18},{\sy+0.38*\hk}) rectangle ++(0.36,{0.62*\hk});
    }
  }
  \fill[sourceblue!80!black] ({4*\wk-0.18},{-3.2+0.38*\hk}) rectangle ++(0.36,{0.62*\hk});
  \fill[sourceblue!80!black] ({1*\wk-0.18},{-6.4+0.38*\hk}) rectangle ++(0.36,{0.62*\hk});
  \fill[sourceblue!80!black] ({5*\wk-0.18},{-6.4+0.38*\hk}) rectangle ++(0.36,{0.62*\hk});
  \node[font=\small, anchor=south] at ({4*\wk},{-3.2+\hk+0.02}) {F$\sharp$};
  \node[font=\small, anchor=south] at ({1*\wk},{-6.4+\hk+0.02}) {D$\flat$};
  \node[font=\small, anchor=south] at ({5*\wk},{-6.4+\hk+0.02}) {A$\flat$};
\end{tikzpicture}
\caption{Three shapes cover nine major chords. Each keyboard shows one example: C major, D major and D\fl{} major.}
\label{fig:three-chord-shapes}
\end{figure}

That leaves three major chords, and each is easy to remember on its own:
\begin{itemize}
  \item G\fl{} major (also called F\sh{} major) is all black keys: G\fl, B\fl, D\fl{}. \listen{Gb3+Bb3+Db4:2}
  \item B\fl{} major is black, white, white: B\fl, D, F. \listen{Bb3+D4+F4:2}
  \item B major is white, black, black: B, D\sh, F\sh{}. \listen{B3+D#4+F#4:2}
\end{itemize}

\begin{intuitionbox}
Minor chords swap the first two shapes. Lower the middle note of C, F or G major and you get white, black, white: Cm is C, E\fl, G. Lower the middle note of D, E or A major and you get three white keys: Dm, Em and Am. So for these six roots, one look tells you the shape of both the major and the minor chord.
\end{intuitionbox}

\section{The chords that belong to a key}

Recall from Chapter~2 that a song ``in the key of C'' uses mostly the notes of the C major scale, the white keys, and feels at home on C. Which chords fit with such a song? The ones built only from those same notes.

To find them, put fingers 1, 3 and 5 on C, E and G. Keep the ``play, skip, play, skip, play'' shape and move your whole hand up one white key at a time. You get seven chords, one on each note of the scale. \listen{C4+E4+G4 D4+F4+A4 E4+G4+B4 F4+A4+C5 G4+B4+D5 A4+C5+E5 B4+D5+F5 C5+E5+G5:2} These are the chords of the key of C.

Musicians number these chords by the scale note they start on, using Roman numerals: I is 1, II is 2, III is 3, IV is 4, V is 5, VI is 6 and VII is 7. A capital numeral means a major chord and a small one means a minor chord. The little circle on vii\textdegree{} means diminished.

\begin{center}
\begin{tabular}{lllll}
\toprule
Numeral & Built on & Chord & Notes & Sound \\
\midrule
I & 1st note & C & C, E, G & major \\
ii & 2nd note & Dm & D, F, A & minor \\
iii & 3rd note & Em & E, G, B & minor \\
IV & 4th note & F & F, A, C & major \\
V & 5th note & G & G, B, D & major \\
vi & 6th note & Am & A, C, E & minor \\
vii\textdegree{} & 7th note & Bdim & B, D, F & diminished \\
\bottomrule
\end{tabular}
\end{center}

Why use numbers instead of names? Every major scale has the same pattern of whole and half steps (Chapter~2). So every major key gives the same pattern of chords: major on 1, 4 and 5, minor on 2, 3 and 6, and diminished on 7.

For example, the key of G uses the G major scale: G, A, B, C, D, E, F\sh. Its chords are G, Am, Bm, C, D, Em and F\sh{}dim. \listen{G3+B3+D4 A3+C4+E4 B3+D4+F#4 C4+E4+G4 D4+F#4+A4 E4+G4+B4 F#4+A4+C5 G4+B4+D5:2} The numbers stay the same: in G, the V chord is D. A song described as ``I, IV, V'' can be played in any key once you know that key's chords.

\subsection*{Home, going away, and coming home}

The chords of a key have different jobs, a bit like places in a story:
\begin{itemize}
  \item I is home. It sounds settled. Most songs end on it.
  \item IV and ii step away from home. They give a feeling of setting out.
  \item V is tension. It pulls strongly back toward I.
\end{itemize}

Hear the pull of V. In C, play C, F, G and stop: the music feels unfinished. \listen{C4+E4+G4:2 F4+A4+C5:2 G3+B3+D4:4} Now let it come home to C. \listen{C4+E4+G4:2 F4+A4+C5:2 G3+B3+D4:2 C4+E4+G4:4}

Part of the pull comes from one note. The G chord contains B, which is only one key below C. Your ear expects that B to rise to C, and when it does, you feel arrival. The vi chord, Am, is home's quieter cousin: it shares two notes, C and E, with the C chord.

\section{Famous chord progressions}

A \textbf{chord progression} is a chain of chords that a song repeats, usually with a steady number of counts on each chord. A few progressions turn up in thousands of songs. Here are four, in the key of C. Each button plays the chords once, two counts each, and then lands on the first chord again.
\begin{itemize}
  \item I--V--vi--IV: C, G, Am, F. One of the most common progressions in pop music. \listen{C4+E4+G4:2 G3+B3+D4:2 A3+C4+E4:2 F4+A4+C5:2 C4+E4+G4:4}
  \item vi--IV--I--V: Am, F, C, G. The same four chords, starting on the minor one, so it feels more wistful. \listen{A3+C4+E4:2 F4+A4+C5:2 C4+E4+G4:2 G3+B3+D4:2 A3+C4+E4:4}
  \item I--vi--IV--V: C, Am, F, G. The sound of many 1950s songs. \listen{C4+E4+G4:2 A3+C4+E4:2 F4+A4+C5:2 G3+B3+D4:2 C4+E4+G4:4}
  \item ii--V--I: Dm, G, C. A short trip home, used everywhere in jazz. \listen{D4+F4+A4:2 G3+B3+D4:2 C4+E4+G4:4}
\end{itemize}

To move a progression to another key, keep the numerals and read the chords from the new key. In the key of G, I is G, ii is Am, IV is C, V is D and vi is Em:

\begin{center}
\begin{tabular}{lll}
\toprule
Numerals & In C & In G \\
\midrule
I--V--vi--IV & C, G, Am, F & G, D, Em, C \\
vi--IV--I--V & Am, F, C, G & Em, C, G, D \\
I--vi--IV--V & C, Am, F, G & G, Em, C, D \\
ii--V--I & Dm, G, C & Am, D, G \\
\bottomrule
\end{tabular}
\end{center}

Hear I--V--vi--IV in G \listen{G3+B3+D4:2 D4+F#4+A4:2 E4+G4+B4:2 C4+E4+G4:2 G3+B3+D4:4} and ii--V--I in G. \listen{A3+C4+E4:2 D4+F#4+A4:2 G3+B3+D4:4}

\begin{tryitbox}
Play I--V--vi--IV in C with your right hand: C, G, Am, F. Count 1, 2, 3, 4 slowly and steadily, and change chord each time you get back to 1. Your hand will jump around for now. In Chapter~4 you will learn to make these changes with almost no movement.
\end{tryitbox}

\section{Seventh chords: adding a fourth note}

Keep going with the ``play, skip, play'' pattern one more step. From C you get C, E, G and then B. A four-note chord like this is a \textbf{seventh chord}, because its top note is seven letters above the root: C, D, E, F, G, A, B.

There are three seventh chords to know first. Their symbols tell you which one you have:
\begin{itemize}
  \item Cmaj7 (``C major seven'') is C, E, G, B. It sounds soft and dreamy. \listen{C4+E4+G4+B4:3}
  \item Dm7 (``D minor seven'') is D, F, A, C. It sounds mellow. \listen{D4+F4+A4+C5:3}
  \item G7 (``G seven'') is G, B, D, F. It has a strong pull toward C. \listen{G3+B3+D4+F4:2 C4+E4+G4:3} Inside it, B wants to rise one key to C, and F wants to fall one key to E.
\end{itemize}

\begin{trapbox}
C7 and Cmaj7 are different chords. C7 is C, E, G, B\fl: its top note is the black key B\fl. Cmaj7 is C, E, G, B. A plain 7 always means the lower top note, which sounds bluesy and restless; maj7 means the higher one. Compare them. \listen{C4+E4+G4+B4:2 C4+E4+G4+Bb4:2}
\end{trapbox}

Four notes in one hand can be a stretch. Here is a trick: play the root low with your left hand, and the other three notes with your right hand. Those top three notes always form a triad you already know:
\begin{itemize}
  \item Cmaj7 is C under E, G, B. The right hand plays Em. \listen{C3+E4+G4+B4:3}
  \item Dm7 is D under F, A, C. The right hand plays F major. \listen{D3+F4+A4+C5:3}
  \item G7 is G under B, D, F. The right hand plays Bdim. \listen{G2+B3+D4+F4:3}
  \item Am7 (A, C, E, G) is A under C, E, G. The right hand plays C major. \listen{A2+C4+E4+G4:3}
\end{itemize}

\begin{labbox}
Play ii--V--I in C with seventh chords: Dm7, G7, Cmaj7. Count 4 slow counts on each chord. Left hand: D3, then G2, then C3, one finger at a time. Right hand: F, A, C, then F, G, B, then E, G, B, all a little above middle C. (In G7 the right hand leaves out D; you will not miss it.) Your right hand barely moves. \listen{@80 D3+F4+A4+C5:4 G2+F4+G4+B4:4 C3+E4+G4+B4:4}
\end{labbox}

\begin{curiousbox}
Place the twelve notes around a clock face, C at 12 o'clock, one note per hour: C\sh{} at 1, D at 2, and so on. Join the three notes of a major chord and you get a triangle whose sides span 4, 3 and 5 hours. Every major chord is the same triangle, turned to a new position. That is why one counting recipe works from every root.

A minor chord is the same triangle flipped over: 3, 4 and 5 hours. A diminished chord is 3, 3 and 6. The chord C, E, G\sh{} (4, 4 and 4) makes a perfect equilateral triangle; it is called augmented.
\end{curiousbox}

\begin{summarybox}
\begin{itemize}
  \item A \emph{melody} is notes one after another; a \emph{chord} is notes played together; \emph{harmony} is the sound of chords under a melody.
  \item A \emph{triad} has three notes: the \emph{root} (which names it), the \emph{third} and the \emph{fifth}. Play it with fingers 1, 3 and 5.
  \item Major: count up 4 keys, then 3 more. Minor: 3, then 4. Only the middle note differs.
  \item Major sounds bright; minor sounds darker. Sus chords sound open and unfinished; diminished chords sound tense.
  \item The chords of a major key, numbered by Roman numerals, are always I, ii, iii, IV, V, vi, vii\textdegree. In C: C, Dm, Em, F, G, Am, Bdim.
  \item I is home, IV steps away, and V pulls back home. A \emph{chord progression} is a chain of chords a song repeats, such as I--V--vi--IV.
  \item A \emph{seventh chord} adds a fourth note. G7 is not Gmaj7. Play the root in the left hand and a familiar triad in the right.
\end{itemize}
\end{summarybox}

\section*{Exercises}

\begin{exercise}
Build E major and E\fl{} major by counting 4 then 3. Which of the three shapes does each belong to? Then lower the middle note of each to get Em and E\fl m.
\end{exercise}

\begin{answerbox}
From E, count 4 keys: F, F\sh, G, G\sh. Then 3 more: A, A\sh, B. So E major is E, G\sh, B. \listen{E4+G#4+B4:2} That is white, black, white.

From E\fl, count 4 keys: E, F, G\fl, G. Then 3 more: A\fl, A, B\fl. So E\fl{} major is E\fl, G, B\fl{}. \listen{Eb4+G4+Bb4:2} That is black, white, black.

Lowering the middle notes gives Em: E, G, B, \listen{E4+G4+B4:2} all white keys. And E\fl m: E\fl, G\fl, B\fl{}, \listen{Eb4+Gb4+Bb4:2} all black keys. We call the middle note G\fl{} rather than F\sh{} so that the letters skip one each time: E, G, B.
\end{answerbox}

\begin{exercise}
Spell Dsus4 (5 then 2 from D) and Dsus2 (2 then 5 from D). Compare them with D major.
\end{exercise}

\begin{answerbox}
Dsus4: from D, count 5 keys to G, then 2 more to A. So Dsus4 is D, G, A. \listen{D4+G4+A4:2} Dsus2: count 2 keys to E, then 5 more to A. So Dsus2 is D, E, A. \listen{D4+E4+A4:2}

D major is D, F\sh, A. All three share D and A. Only the middle note changes: G for sus4, F\sh{} for major, E for sus2. Hear Dsus4 settle into D major. \listen{D4+G4+A4:2 D4+F#4+A4:3}
\end{answerbox}

\begin{exercise}
The F major scale is F, G, A, B\fl, C, D, E. List the seven chords of the key of F, with their Roman numerals. Then play I--V--vi--IV in F.
\end{exercise}

\begin{answerbox}
Use the ``play, skip, play, skip, play'' shape on the notes of the F major scale. You get: I is F (F, A, C); ii is Gm (G, B\fl, D); iii is Am (A, C, E); IV is B\fl{} (B\fl, D, F); V is C (C, E, G); vi is Dm (D, F, A); vii\textdegree{} is Edim (E, G, B\fl).

So I--V--vi--IV in F is F, C, Dm, B\fl{}. \listen{F3+A3+C4:2 C4+E4+G4:2 D4+F4+A4:2 Bb3+D4+F4:2 F3+A3+C4:4}
\end{answerbox}

\begin{exercise}
A song in the key of G uses only three chords: G, C and D. What are their Roman numerals? Which three chords would the same song use in the key of D? (The D major scale is D, E, F\sh, G, A, B, C\sh.)
\end{exercise}

\begin{answerbox}
In G, the chord on the 1st note is G, the 4th note is C, and the 5th note is D. So the song uses I, IV and V. In the key of D, the 1st, 4th and 5th notes are D, G and A, so the song uses D, G and A major. \listen{D4+F#4+A4:2 G3+B3+D4:2 A3+C#4+E4:2 D4+F#4+A4:4}
\end{answerbox}

\begin{exercise}
Em7 is E, G, B, D. Which triad does the right hand play over a left-hand E? And D7 is D, F\sh, A, C: which two of its notes move by one key to reach the G chord (G, B, D)?
\end{exercise}

\begin{answerbox}
Take away the root E and you are left with G, B, D: a G major triad. So Em7 is E in the left hand under G major in the right hand. \listen{E3+G4+B4+D5:3}

In D7, F\sh{} rises one key to G, and C falls one key to B. That is the same pull you met in G7 going to C, moved to a new key. \listen{D4+F#4+A4+C5:2 D4+G4+B4:3}
\end{answerbox}
